\documentclass[10pt,a4paper]{amsart}
\usepackage[margin=19mm]{geometry}
\usepackage{amsmath,amssymb,mathtools}
\usepackage[scr=rsfs,cal=euler]{mathalfa}
\usepackage{xfrac}
\usepackage[unicode,hidelinks,bookmarks=false]{hyperref}
\newcommand{\ie}{i.e.\ }
\newcommand{\cf}{cf.\ }
\newcommand{\resp}{respectively\ }
\newcommand{\Subsection}{\subsection}
\providecommand{\nobreakdash}{\nobreak\hbox{-}}
\setlength{\emergencystretch}{2em}
\multlinegap0pt
\usepackage{bm}
\usepackage{bbm}
\usepackage{dsfont}
\usepackage{xspace}
\usepackage{cancel}
\usepackage{mathtools}
\usepackage{amsthm}
\makeatletter
\@ifundefined{theorem}{\newtheorem{theorem}{Theorem}}{}
\@ifundefined{lemma}{\newtheorem{lemma}[theorem]{Lemma}}{}
\makeatother
\def\deg{\mathrm{deg}}
\let\rho=\varrho 
\newcommand{\Ptilde}{\widetilde P}
\title[On the self-intersection time of non-backtracking random wal]{On the self-intersection time of non-backtracking random walks}
\author{Ferenc Bencs, Leslie Ann Goldberg, Matthew Jenssen, Mark Jerrum, Gabor Pete, Guus Regts, Yitong Yin}
\date{Source: arXiv:2608.09729v1 (CC BY 4.0)}
\begin{document}
\maketitle

\noindent\emph{Source and licence.} On the self-intersection time of non-backtracking random walks. Version arXiv:2608.09729v1 (CC BY 4.0), distributed under the Creative Commons Attribution 4.0 International licence, as recorded in the arXiv licence metadata for this exact version and shown on the abstract page https://arxiv.org/abs/2608.09729v1. Full rights notice: licenses/arxiv-2608.09729-CC-BY-4.0.txt.\par\smallskip

\noindent\textit{Editorial scope.} Counted unit: the Lemma whose source label is \texttt{lem:endpoint-mult} in the article (source0.tex lines 234--241, source label \texttt{lem:endpoint-mult}), delivered together with the article's own complete original proof of it (source0.tex lines 243--350, one \texttt{proof} environment of 108 lines). No mathematics is written, repaired or re-derived: statement and proof are the article's own text line for line. Required context retained uncounted: none. Every definition, display and label of those lines is retained unchanged. Omitted from this package and not redistributed: 1--233, 242--242, 351--1286. Nothing inside a retained excerpt is omitted or altered: every retained body is delivered exactly as the article's lines are written, so no omission receipt of any kind accompanies this package. Citations: the retained excerpts carry no citation command at all, so no bibliography entry of the article is transcribed and no reference is redistributed. Reference adaptation: none was needed and none was made; every reference inside the delivered unit resolves inside it, so no unresolved reference or citation is produced. Printed slips kept literal: no printed word, symbol, hypothesis or number of the counted statement or of its proof was altered. Layout and class: the article's own class is \texttt{article} with packages including a4wide, amsmath, amsthm, tikz, amsfonts, amssymb, color, thm-restate, comment; the wrapper is the lane's pinned \texttt{amsart} editorial wrapper at 10pt on A4, which restates the theorem-like environments the retained text needs and adds the article's own macro definitions that the retained text needs (\texttt{\textbackslash deg}, \texttt{\textbackslash rho}), with extra wrapper packages (bm, bbm, dsfont, xspace, cancel, mathtools). The delivered numbering is the unit's own. Assets: no retained span contains a figure environment, a graphics-inclusion command, a PDF-page inclusion, a table or a TikZ picture; no image file is copied into this package, the package declares no asset dependency and no image file is redistributed with it. Licence: version arXiv:2608.09729v1 of this article is distributed under the Creative Commons Attribution 4.0 International licence, as recorded in the arXiv licence metadata for this exact version and shown on the abstract page https://arxiv.org/abs/2608.09729v1. A full rights notice is delivered at licenses/arxiv-2608.09729-CC-BY-4.0.txt. Characters: every retained excerpt is pure ASCII, so the delivered engine, XeLaTeX with its default Latin Modern text font, reports no missing character.
\begin{lemma}\label{lem:endpoint-mult}
If $P=v_0\cdots v_m$ is an NBW, and $w\in V(G)$, then
\[
|\mathcal A_w(P)|
\le
\lceil\log_2 d\rceil.
\]
\end{lemma}

\begin{proof}
 Suppose that
$i\in F(P)$ and $\rho\in L'(i,P)$. By the definition of
$L'(i,P)$, there is a leaf $\tilde \rho\in L(i,P)$ that is a
descendant of $\rho$. Consequently, for every prefix $Q$ of
$\ell(\rho)$, and every node
$\eta\in T_{Q_{-2},Q_{-1}}$, we have
\[
P_{\le i}\circ Q\circ\ell(\eta)\notin\mathcal I.
\tag{1}\label{eq:safe-descendant}
\]
Indeed, $Q$ is also a prefix of $\ell(\tilde\rho)$, so this follows
directly from the definition of $L(i,P)$.

Consider now
\[
(i,\rho),(j,\rho')\in\mathcal A_w(P), \quad i\leq j\, .
\]
We will show that $\ell(\rho)$ is a suffix of $\ell(\rho')$ or vice versa. 
Suppose, for a contradiction, that neither $\ell(\rho)$ nor
$\ell(\rho')$ is a suffix of the other. Since both walks end at $w$,
they have a nonempty maximal common suffix. Let $u$ be the first
vertex of this maximal common suffix. By assumption, the predecessors of $u$ in
$\ell(\rho)$ and $\ell(\rho')$ are distinct; denote them by $x$ and
$y$ respectively.

Let
\[
A=a_0a_1\cdots a_r
\]
be the prefix of $\ell(\rho)$ ending at $a_r=u$, so that
$a_{r-1}=x$, and let $Q$ be the prefix of $\ell(\rho')$ ending at
$u$. Define
\[
R:=yua_{r-1}a_{r-2}\cdots a_0.
\]
Because $x\neq y$ and the reverse of an NBW is an NBW, $R$ is an
NBW beginning with the last directed edge $yu$ of $Q$.

The mass of $R$ in $T_{y,u}$ satisfies
\[
\pi(R)
=
\frac{\pi_i(A)}{\deg(u)-1}
\ge
\frac{\pi_i(\rho)}{d}
\ge
\frac{1}{\sqrt n}\, .
\]
Here we used that $A$ is a prefix of $\ell(\rho)$, that
$\deg(u)-1\le d$, and that
$\pi_i(\rho)\ge d/\sqrt n$. It follows that $R$ is the label of
a node of $T_{y,u}$.

On the other hand,
\[
P_{\le j}\circ Q\circ R\in\mathcal I,
\]
because $R$ returns to $a_0=v_{i-1}$, which has already appeared
in $P_{\le j}$. This contradicts
\eqref{eq:safe-descendant}, applied to $\rho'$ and the prefix
$Q$. Thus $\ell(\rho)$ is a suffix of $\ell(\rho')$ or vice versa as claimed. 

We also note that if 
\(
(i,\rho),(j,\rho')\in\mathcal A_w(P)
\)
are distinct then the labels $\ell(\rho), \ell(\rho')$ are
distinct.
Indeed, since $j\in F(P)$, $P_{\le j}$ is
self-avoiding: otherwise $L(j,P)=\emptyset$, contradicting
$j\in F(P)$. Hence, if
$\ell(\rho)=\ell(\rho')$, equality of their first directed edges
implies $i=j$ ($v_i=v_j$ implies $i=j$ since $P_{\leq j}$ is self-avoiding). Since nodes of a fixed tree are uniquely determined
by their labels, this then implies $\rho=\rho'$.

Now choose $(i_0,\rho_0)\in\mathcal A_w(P)$ so that
$\ell(\rho_0)$ has minimum length. By the above, we have that $\ell(\rho_0)$ is a suffix of every other label
in $\mathcal A_w(P)$.

For $(i,\rho)\in\mathcal A_w(P)$, let $q=q(i,\rho)$ be the
number of edges of $\ell(\rho)$ preceding 
$\ell(\rho_0)$. Every transition probability along these $q$ edges
is at most $1/2$, because the minimum degree of $G$ is at least $3$.
Therefore
\[
\pi_i(\rho)
\le
2^{-q}\pi_{i_0}(\rho_0).
\]
Using the defining mass bounds for $L'(i,P)$ and
$L'(i_0,P)$, we obtain
\[
\frac{d}{\sqrt n}
\le
\pi_i(\rho)
\le
2^{-q}\pi_{i_0}(\rho_0)
<
2^{-q}\bigl(\deg(w)-1\bigr)\frac{d}{\sqrt n}.
\]
It follows that
\[
0\le q<
\log_2\bigl(\deg(w)-1\bigr).
\]
The result follows by recalling that each element of $\mathcal{A}_w(P)$ has a unique label and distinct labels give distinct values of $q$ (since all labels are suffix-comparable). The result follows. 
\end{proof}

\end{document}
